%% The following is a directive for TeXShop to indicate the main file
%%!TEX root = diss.tex

%=======================================
\chapter{Operators}
\label{ch:Operators}
%=======================================

In this section we'll make our way towards the larger problem of higher-dimensional labelled walks. The methods will be similar to those in the section on the one-dimensional problem but their complexity will rise substantially. Without the operator notation even the derivation of the functional equation is quite laborious. It's probably worth mentioning that we found the derivation the hard way first before we discovered we can use operators to simplify the process. What once took about a dozen pages now takes a just few. In this section we'll explore the properties of the operators that interest us, showing when the commute, if they have inverses, what their eigenvalues and eigenvectors are and so on.

%=======================================
\section{Definitions}
\label{sec:Definitions}
%=======================================

We'll work with many dimensions directly as that doesn't really change the situation very much. Let $d$ be the number of dimensions. A walk in $d$ dimensions is valid with respect to coordinate $x_i$ if it never steps on a point with negative $x_i$ coordinate. A walk in $d$ dimensions is valid if it is valid with respect to all $x_i$ coordinates, that is it remains in the non-negative orthant. Let
\[
	E_i = 1 - \ol{x}_i[\ol{x}_i]
\]
be the extraction operator with respect to $x_i$. Applying it on an expression like $SQ$ yields the generating function for $x_i$-valid walks. Similarly let 
\[
	L_i = 1 + b_i [x_i^0]
\]
be the labelling operator with respect to the $x_i = 0$ hyperplane. This operator labels walks that end on the $x_i = 0$ hyperplane with the label $a_i$. Let
\[
	M_i = 1 - c_i[x_i^0].
\]
We can call this operator the de-labelling operator because it is the inverse of $L_i$, which we'll show later. Let
\[
	B_i = 1 - [b_i^0].
\]
We don't have a name for this operator yet but we use it when trying to establish a relationship between labelled and unlabelled quantities. Maybe we can call it the relator? Next let
\[
	P = P_d \ldots P_1
\]
For any operator $P$, omitting the subscript will mean
\[
	P = P_d\ldots P_1
\]

%=======================================
\section{Commutativity properties}
\label{sec:CommProp}
%=======================================

Firstly we have that if $i \neq 0$
\[
	[x^i][x^k] = 0
\]
Then
\[
	[x^i]x^j[x^k] = \delta_{i,j}[x^k]
\]
This helps establish the commutativity properties of all the operators we have so far $E_i, L_i, M_i, D_i, B_i$. In general, if we have operators $A$ and $B$ given by
\[
	A(x) = \sum_k A_k x^k[x^k],\quad B(x) = \sum_k B_k x^k[x^k]
\]
then
\begin{align*}
	A(x)B(x) &= \left(\sum_k A_k x^k[x^k]\right)\left(\sum_k B_k x^k[x^k]\right)\\
	& = \sum_{i,j}A_ix^i[x^i]B_jx^j[x^j] \\
	&= \sum_{i,j}A_iB_j x^i\delta_{i,j}[x^j]\\
	&= \sum_k A_kB_k x^k[x^k]
\end{align*}
We obtain the same result if we reverse $A$ and $B$, so they commute. Now
\begin{align*}
	A(x)B(y) &= \left(\sum_k A_k x^k[x^k]\right)\left(\sum_k B_k y^k[y^k]\right)\\
	& = \sum_{i,j}A_ix^i[x^i]B_jy^j[y^j]\\
	& = \sum_{i,j}A_iB_jx^iy^j[x^iy^j]\\
	& = \sum_{i,j}B_jy^j[y^j]A_ix^i[x^i]\\
	& = \left(\sum_k B_k y^k[y^k]\right)\left(\sum_k A_k x^k[x^k]\right)\\
	&  = B(y)A(x)
\end{align*}
So they commute. Our operators are of the form $1 + A(x)$ so they will commute. In general, as long as each coefficient extraction is only doing things on the level of the power of $x$ they are extracting, the operators will commute. If there is some mixing we'll potentially not have commuting operators. This allows us to later create easy generalizations of the labelling problem so far.

%=======================================
\section{Kernels, eigenvalues and eigenfunctions}
\label{sec:KernelStuff}
%=======================================

%=======================================
\subsection{Operator $L$}
\label{sec:OpL}
%=======================================

Operator $L$ is invertible with inverse $M$
\begin{align*}
	L M &= (1+b[x^0])(1-c[x^0]) \\
	& = 1+b[x^0]-c[x^0]-b[x^0]c[x^0]\\
	& = 1 + (b-c-bc)[x^0]\\
	& = 1 + (a-1-1+\ol{a}-a+1+1-\ol{a})[x^0]\\
	& = 1.
\end{align*}
As a bonus we see that it has a trivial kernel. Next, we'll determine its eigenvalues. We expand the equation $Lf = \lambda f$ to get
\begin{align*}
	Lf &= f + b f_0 = \lambda f\\
	(1-\lambda)f &= -bf_0\\
	f&=\frac{b f_0}{\lambda - 1}
\end{align*}
Therefore we determine that $f$ is a constant with respect to $x$. So the operator's eigenvalues are all real functions $\lambda(a,t) \neq 1$ the eigenfunctions are $f(x,a,t) = \alpha/(\lambda - 1)$. This is admitedly not a very interesting situation. What if $\lambda = 1$. Then any function $f$ that is non-zero at $x = 0$ will be an eigenfunction, that is functions of the form $f(x,a,t) = x g(x,a,t)$ where $g$ is finite at $x = 0$.

%=======================================
\subsection{Operator $E$}
\label{sec:OpE}
%=======================================

We take $E = 1 - \ol{x} [\ol{x}]$. We have
\[
	0 = E(f) = f - \ol{x}f_{-1} = 0
\]
which means that $f$ is in the kernel of $E$ if $f = c\ol{x}$ for some constant. This shows that $E$ itself is not invertible with kernel generated by $\ol{x}$. Next we determine the eigenvalues
\begin{align*}
	Ef &= f - \ol{x} f_{-1} = \lambda f\\
	(1-\lambda)f &= f_{-1}\ol{x}\\
	f&=\frac{f_{-1}}{x(\lambda - 1)}
\end{align*}
If $f_{-1} = 0$ then $E$ acts as the identity operator so the eigenvalue is $\lambda = 1$ for every such function. If $f_{-1} \neq 0$ then
\[
	Ef = \frac{f_{-1}}{x(\lambda - 1)} - \frac{1}{x}\frac{f_{-1}}{(\lambda - 1)} = 0
\]
so the eigenvalue must be 0. Therefore there is a splitting between the functions $E$ acts on. Those without an $\ol{x}$ are preserved with eigenvalue 1. Those who have an $\ol{x}$ and other terms are not preserved. Finally, if $f$ is just such a term, we have eigenvalue 0, of course.

%=======================================
\subsection{Operator $B$}
\label{sec:OpB}
%=======================================

The situation here is analogous to the one for $L$ with a different variable and coefficient $-1$.

%=======================================
\subsection{Compounded operators}
\label{sec:CompoundOps}
%=======================================

Our first example concerns $ES$. Since the kernel of $E$ is generated by $\ol{x}$ then a function $f$ is in the kernel of $ES$ if and only if $fS$ is in the kernel of $E$. That is, the kernel of $ES$ is generated by
\[
	f = \frac{c}{xS} = \frac{c}{S_{-1}+S_0 x + S_1 x^2}
\]
This function can have various forms depending on which coefficients are zero.

Finally we should look at the invertibility and other properties of the operator $1-tLES$ or even $M-tES$. Okay this is difficult. I'll have to spend more time on this later.
%=======================================
\subsection{Operator $\mc{L} = 1- tLES$}
\label{sec:OpEqn}
%=======================================
\endinput
